% fgib_linear.tex —— FGIB 旁路瓶颈与 CEB 主路瓶颈的等价性推导
% 记号沿用 Fischer (2020) "The Conditional Entropy Bottleneck"
% 用法：导言区 \usepackage{amsmath,amssymb,amsthm}，正文处 % fgib_linear.tex —— FGIB 旁路瓶颈与 CEB 主路瓶颈的等价性推导
% 记号沿用 Fischer (2020) "The Conditional Entropy Bottleneck"
% 用法：导言区 \usepackage{amsmath,amssymb,amsthm}，正文处 % fgib_linear.tex —— FGIB 旁路瓶颈与 CEB 主路瓶颈的等价性推导
% 记号沿用 Fischer (2020) "The Conditional Entropy Bottleneck"
% 用法：导言区 \usepackage{amsmath,amssymb,amsthm}，正文处 % fgib_linear.tex —— FGIB 旁路瓶颈与 CEB 主路瓶颈的等价性推导
% 记号沿用 Fischer (2020) "The Conditional Entropy Bottleneck"
% 用法：导言区 \usepackage{amsmath,amssymb,amsthm}，正文处 \input{fgib_linear.tex}
% 若论文已定义 lemma/proposition/remark 环境，删除下面三行 \newtheorem

\newtheorem{lemma}{Lemma}
\newtheorem{proposition}{Proposition}
\newtheorem{remark}{Remark}

\textbf{Notation.}
Let $X$ be the input and $Y$ the label.
The main (deterministic) path of FGIB computes the representation
$H = f_{\mathrm{enc}}(X) \in \mathbb{R}^{d}$, and the classifier consumes $H$ directly.
The side path parameterizes $q(z|x) = \mathcal{N}(\mu_q(x), \Sigma_q(x))$ with linear mean
head $\mu_q(x) = A\,h(x)$ for $A \in \mathbb{R}^{z \times d}$, unconstrained ($z = d$, so
$A$ is a.s.\ invertible; the bias of the implementation's Linear head is omitted by the
same convention); $Z$ never feeds the classifier.
The label-conditional prior (backward encoder) is $r(z|y) = \mathcal{N}(\mu_p(y), \Sigma_p(y))$;
for FGIB it is the fixed orthonormal anchor table $\mathcal{N}(a v_y, I)$. The
orthogonality of the method lies in this fixed frame and in the equilibrium it induces
--- the KL drives the class-conditional means $\mu_y = A\,\mathbb{E}[h|y]$ toward the
orthonormal anchors $a v_y$, so the transformed class representations become mutually
orthogonal (the aligned family of fgib\_degeneration.tex) --- while $A$ itself remains
an unconstrained linear map.
The side-path regularizer is $\operatorname{KL}(q(z|x)\,\|\,r(z|y))$, added to the
cross-entropy with weight $\beta$, with gradients reaching the main path only through the
shared encoder (the classification gradient does not enter the side heads, and vice versa).
Following Fischer (2020), we write the residual information $I(X;Z|Y)$ (Eq.~4), the CEB
objective $\min_Z I(X;Z|Y) - \gamma I(Y;Z)$ (Eq.~5), the MNI point
$I(X;Y) = I(X;Z) = I(Y;Z)$ (Eq.~1), and the variational upper bound
$I(X;Z|Y) \le \mathbb{E}[\operatorname{KL}(q(z|x)\,\|\,r(z|y))]$ (Eqs.~9--11).

\textbf{Motivation.}
Stochastic bottlenecks such as VIB and CEB are \emph{stochastic at training time but
deterministic at evaluation time}: the classifier is trained on sampled codes $z \sim q(z|x)$
yet evaluated on the mean $\mu_q(x)$ (Fischer 2020, Appendix A.1), so its training and
evaluation distributions differ. In the same models the cross-entropy gradient and the KL
gradient act on the same heads: the classification term can exploit the variance (by Jensen,
$\mathbb{E}_z[\mathrm{CE}(c(z), y)] \ge \mathrm{CE}(c(\mu_q), y)$, so loosening the noise
relaxes the upper bound), while the KL term pushes the posterior back toward the prior.
Finally, the regularizer constrains the noisy code $z$, whereas inference consumes the
noise-free mean --- the object being regularized is not the object being used.
FGIB resolves all three mismatches by design: the classifier consumes the deterministic
representation $h$ at both training and evaluation time, the two objectives act on disjoint
parameter sets (the classification gradient does not enter the side heads, and the KL
gradient reaches the main path only through the shared encoder), and, by
Proposition~\ref{prop:fgib-ceb} below, the side-path regularizer is exactly a CEB bottleneck
applied to the representation that inference actually uses. Confining the stochasticity to a
side path also keeps the variational objective well-defined: for continuous deterministic
representations the residual information diverges (Amjad \& Geiger; Kolchinsky et al.),
so a purely deterministic bottleneck admits no finite variational bound, while the finite
side-path covariance of FGIB does (Remark~\ref{rem:limit}). Whether the robustness benefits
of the stochastic bottleneck survive the deterministic inference path is the empirical
question studied in the experiments.

\begin{lemma}[Linear invariance of the Gaussian KL]\label{lem:kl-invariance}
Let $q = \mathcal{N}(\mu_q, \Sigma_q)$ and $p = \mathcal{N}(\mu_p, \Sigma_p)$ on $\mathbb{R}^{z}$,
and let $A \in \mathbb{R}^{z \times z}$ be invertible. Define the common pushforwards
$q' = A_\# q = \mathcal{N}(A\mu_q, A\Sigma_q A^\top)$ and $p' = A_\# p$. Then
$\operatorname{KL}(q'\,\|\,p') = \operatorname{KL}(q\,\|\,p)$.
\end{lemma}

\begin{proof}
With $d = z$, the closed form reads
$\operatorname{KL}(q\,\|\,p)
= \tfrac{1}{2}\big[\operatorname{tr}(\Sigma_p^{-1}\Sigma_q)
+ (\mu_q - \mu_p)^\top \Sigma_p^{-1}(\mu_q - \mu_p)
- d + \ln|\Sigma_p| - \ln|\Sigma_q|\big]$.
Under the common map $A$ the three non-constant terms are unchanged:
\begin{align*}
\operatorname{tr}\big((A\Sigma_p A^\top)^{-1}(A\Sigma_q A^\top)\big)
&= \operatorname{tr}(A^{-\top}\Sigma_p^{-1}A^{-1} A\,\Sigma_q A^\top)
 = \operatorname{tr}(A^{-\top}\Sigma_p^{-1}\Sigma_q A^\top) \\
&= \operatorname{tr}(\Sigma_p^{-1}\Sigma_q A^\top A^{-\top})
 = \operatorname{tr}(\Sigma_p^{-1}\Sigma_q), \\[2pt]
(A\Delta\mu)^\top (A\Sigma_p A^\top)^{-1} (A\Delta\mu)
&= \Delta\mu^\top A^\top A^{-\top}\Sigma_p^{-1}A^{-1}A\,\Delta\mu
 = \Delta\mu^\top \Sigma_p^{-1}\Delta\mu, \\[2pt]
\ln|A\Sigma_p A^\top| - \ln|A\Sigma_q A^\top|
&= \big(\ln|\Sigma_p| + 2\ln|A|\big) - \big(\ln|\Sigma_q| + 2\ln|A|\big)
 = \ln|\Sigma_p| - \ln|\Sigma_q|,
\end{align*}
where $\Delta\mu = \mu_q - \mu_p$ and the trace step uses the cyclic property. \hfill\qed
\end{proof}

\begin{lemma}[Invariance of the CEB objective and the MNI point]\label{lem:mi-invariance}
For any bijective measurable map $\phi$,
$I(X;\phi(Z)|Y) = I(X;Z|Y)$ and $I(Y;\phi(Z)) = I(Y;Z)$.
Consequently the CEB objective (Eq.~5) is invariant under bijections of $Z$, and
$Z$ attains the MNI point (Eq.~1) if and only if $\phi(Z)$ does.
\end{lemma}

\begin{proof}
Mutual information is invariant under bijective transformations of either argument
(data processing inequality applied in both directions). \hfill\qed
\end{proof}

\begin{proposition}[The FGIB side path is a CEB bottleneck on the main path]\label{prop:fgib-ceb}
Let $z = d$ and let the learned head $A$ be invertible (unconstrained). The \emph{fixed}
prior $r(z|y) = \mathcal{N}(\mu_p(y), \Sigma_p(y))$ --- for FGIB, the orthonormal anchor
table $\mu_p(y) = a\,v_y$, $\Sigma_p = I$ --- is, by construction, the image under $A$
of the canonical prior
$r'(\tilde{z}|y) = \mathcal{N}(A^{-1}\mu_p(y), A^{-1}\Sigma_p(y) A^{-\top})$
(since $A A^{-1} = I$), and the posterior $q(z|x) = \mathcal{N}(A\,h(x), \Sigma_q(x))$
is the image of
$q'(\tilde{z}|x) = \mathcal{N}(h(x), A^{-1}\Sigma_q(x) A^{-\top})$.
The prior is never transformed by training: $A^{-1}$ only reads the fixed anchor
geometry in the coordinates in which the posterior is centered at $h$.
By KL invariance under the common invertible map $A^{-1}$
(Lemma~\ref{lem:kl-invariance}), pointwise for every $(x, y)$,
\begin{equation}
\operatorname{KL}\big(q(z|x)\,\|\,r(z|y)\big)
= \operatorname{KL}\big(q'(\tilde{z}|x)\,\|\,r'(\tilde{z}|y)\big),
\label{eq:fgib-ceb}
\end{equation}
That is, the side-path regularizer is pointwise identical to a CEB bottleneck applied to
a noisy version of the main-path representation,
$\tilde{z} = h + \tilde{\epsilon}$ with
$\tilde{\epsilon} \sim \mathcal{N}(0, A^{-1}\Sigma_q A^{-\top})$,
using the canonical prior $r'$ as the backward encoder. For FGIB the canonical backward
encoder is the frame $\mathcal{N}(a\,A^{-1} v_y, A^{-1}A^{-\top})$, whose anchors
$A^{-1}v_y$ are orthonormal at scale $a$ in the induced metric $A^{\top}A$
($\langle A^{-1}v_i, A^{-1}v_j\rangle_{A^{\top}A} = \delta_{ij}$): the orthogonality
enters through the fixed frame $\{v_y\}$ and the equilibrium it induces, not through a
constraint on $A$.
Hence the regularizer's gradient with respect to $h$ (through the shared encoder)
coincides, up to the fixed linear map $A^{\top}$, with the gradient that a CEB model
applies to its representation mean.
\end{proposition}

\begin{proof}
The map $z \mapsto \tilde{z} = A^{-1}z$ is a bijection, and
$q' = (A^{-1})_\# q$, $r' = (A^{-1})_\# r$:
$A^{-1}A\,h = h$ and $A^{-1}\Sigma_q A^{-\top}$ give the canonical posterior, and
$A^{-1}\mu_p$, $A^{-1}\Sigma_p A^{-\top}$ give the canonical prior, with
$A\,(A^{-1}\mu_p) = \mu_p$ and $A\,(A^{-1}\Sigma_p A^{-\top})\,A^{\top} = \Sigma_p$, so
$r$ is exactly the image of $r'$ under $A$.
Equation~\eqref{eq:fgib-ceb} is Lemma~\ref{lem:kl-invariance} applied to $A^{-1}$.
Sampling $\tilde{z} \sim q'$ is equivalent to
$\tilde{z} = h + \tilde{\epsilon}$ with the stated noise.
Combined with the variational bound of Fischer (2020, Eqs.~9--11), the regularizer
upper-bounds the residual information of the noisy main-path representation:
$I(X;\tilde{Z}|Y) \le \mathbb{E}[\operatorname{KL}(q'(\tilde{z}|x)\,\|\,r'(\tilde{z}|y))]$.
\hfill\qed
\end{proof}

\begin{remark}[Family closure]\label{rem:closure}
Lemma~\ref{lem:kl-invariance} holds for the underlying Gaussian distributions. Within a
restricted variational family the equality of Proposition~\ref{prop:fgib-ceb} is exact
only if the family is closed under the inverse map $A^{-1}$: the full-covariance family
is closed (and is the parameterization used by Fischer 2020); the diagonal family is not
--- for diagonal $q, r$ the transformed pair is diagonal only if $A^{-1}$ preserves
diagonality (e.g., permutations and sign flips) and approximate otherwise. The diagonal
implementation of the codebase is therefore covered at the distribution level, with the
family approximation confined to the noise covariance $A^{-1}\Sigma_q A^{-\top}$
(fgib\_bound.tex, Remark on the unconstrained side head).
\end{remark}

\begin{remark}[Rank deficiency]\label{rem:rank}
If $z < d$ or $A$ is rank-deficient, only the projection of $h$ onto the image of $A$
receives gradient from the regularizer; the null-space component is unconstrained.
In the default configuration $z = d$ and a randomly initialized $A$ is almost surely invertible.
\end{remark}

\begin{remark}[Noise-free limit]\label{rem:limit}
The equivalence of Proposition~\ref{prop:fgib-ceb} should be understood at the level of the
variational objective, not as a mutual-information limit: as the side-path noise vanishes,
$I(X;Z|Y)$ diverges for continuous deterministic $H$ (the pathology of deterministic
networks discussed by Amjad \& Geiger and Kolchinsky et al.), and the variational KL bound
becomes loose. Keeping the side-path covariance finite --- without injecting noise into the
inference path --- is precisely what keeps the regularizer well-defined; this is the design
rationale of FGIB.
\end{remark}

% 若论文已定义 lemma/proposition/remark 环境，删除下面三行 \newtheorem

\newtheorem{lemma}{Lemma}
\newtheorem{proposition}{Proposition}
\newtheorem{remark}{Remark}

\textbf{Notation.}
Let $X$ be the input and $Y$ the label.
The main (deterministic) path of FGIB computes the representation
$H = f_{\mathrm{enc}}(X) \in \mathbb{R}^{d}$, and the classifier consumes $H$ directly.
The side path parameterizes $q(z|x) = \mathcal{N}(\mu_q(x), \Sigma_q(x))$ with linear mean
head $\mu_q(x) = A\,h(x)$ for $A \in \mathbb{R}^{z \times d}$, unconstrained ($z = d$, so
$A$ is a.s.\ invertible; the bias of the implementation's Linear head is omitted by the
same convention); $Z$ never feeds the classifier.
The label-conditional prior (backward encoder) is $r(z|y) = \mathcal{N}(\mu_p(y), \Sigma_p(y))$;
for FGIB it is the fixed orthonormal anchor table $\mathcal{N}(a v_y, I)$. The
orthogonality of the method lies in this fixed frame and in the equilibrium it induces
--- the KL drives the class-conditional means $\mu_y = A\,\mathbb{E}[h|y]$ toward the
orthonormal anchors $a v_y$, so the transformed class representations become mutually
orthogonal (the aligned family of fgib\_degeneration.tex) --- while $A$ itself remains
an unconstrained linear map.
The side-path regularizer is $\operatorname{KL}(q(z|x)\,\|\,r(z|y))$, added to the
cross-entropy with weight $\beta$, with gradients reaching the main path only through the
shared encoder (the classification gradient does not enter the side heads, and vice versa).
Following Fischer (2020), we write the residual information $I(X;Z|Y)$ (Eq.~4), the CEB
objective $\min_Z I(X;Z|Y) - \gamma I(Y;Z)$ (Eq.~5), the MNI point
$I(X;Y) = I(X;Z) = I(Y;Z)$ (Eq.~1), and the variational upper bound
$I(X;Z|Y) \le \mathbb{E}[\operatorname{KL}(q(z|x)\,\|\,r(z|y))]$ (Eqs.~9--11).

\textbf{Motivation.}
Stochastic bottlenecks such as VIB and CEB are \emph{stochastic at training time but
deterministic at evaluation time}: the classifier is trained on sampled codes $z \sim q(z|x)$
yet evaluated on the mean $\mu_q(x)$ (Fischer 2020, Appendix A.1), so its training and
evaluation distributions differ. In the same models the cross-entropy gradient and the KL
gradient act on the same heads: the classification term can exploit the variance (by Jensen,
$\mathbb{E}_z[\mathrm{CE}(c(z), y)] \ge \mathrm{CE}(c(\mu_q), y)$, so loosening the noise
relaxes the upper bound), while the KL term pushes the posterior back toward the prior.
Finally, the regularizer constrains the noisy code $z$, whereas inference consumes the
noise-free mean --- the object being regularized is not the object being used.
FGIB resolves all three mismatches by design: the classifier consumes the deterministic
representation $h$ at both training and evaluation time, the two objectives act on disjoint
parameter sets (the classification gradient does not enter the side heads, and the KL
gradient reaches the main path only through the shared encoder), and, by
Proposition~\ref{prop:fgib-ceb} below, the side-path regularizer is exactly a CEB bottleneck
applied to the representation that inference actually uses. Confining the stochasticity to a
side path also keeps the variational objective well-defined: for continuous deterministic
representations the residual information diverges (Amjad \& Geiger; Kolchinsky et al.),
so a purely deterministic bottleneck admits no finite variational bound, while the finite
side-path covariance of FGIB does (Remark~\ref{rem:limit}). Whether the robustness benefits
of the stochastic bottleneck survive the deterministic inference path is the empirical
question studied in the experiments.

\begin{lemma}[Linear invariance of the Gaussian KL]\label{lem:kl-invariance}
Let $q = \mathcal{N}(\mu_q, \Sigma_q)$ and $p = \mathcal{N}(\mu_p, \Sigma_p)$ on $\mathbb{R}^{z}$,
and let $A \in \mathbb{R}^{z \times z}$ be invertible. Define the common pushforwards
$q' = A_\# q = \mathcal{N}(A\mu_q, A\Sigma_q A^\top)$ and $p' = A_\# p$. Then
$\operatorname{KL}(q'\,\|\,p') = \operatorname{KL}(q\,\|\,p)$.
\end{lemma}

\begin{proof}
With $d = z$, the closed form reads
$\operatorname{KL}(q\,\|\,p)
= \tfrac{1}{2}\big[\operatorname{tr}(\Sigma_p^{-1}\Sigma_q)
+ (\mu_q - \mu_p)^\top \Sigma_p^{-1}(\mu_q - \mu_p)
- d + \ln|\Sigma_p| - \ln|\Sigma_q|\big]$.
Under the common map $A$ the three non-constant terms are unchanged:
\begin{align*}
\operatorname{tr}\big((A\Sigma_p A^\top)^{-1}(A\Sigma_q A^\top)\big)
&= \operatorname{tr}(A^{-\top}\Sigma_p^{-1}A^{-1} A\,\Sigma_q A^\top)
 = \operatorname{tr}(A^{-\top}\Sigma_p^{-1}\Sigma_q A^\top) \\
&= \operatorname{tr}(\Sigma_p^{-1}\Sigma_q A^\top A^{-\top})
 = \operatorname{tr}(\Sigma_p^{-1}\Sigma_q), \\[2pt]
(A\Delta\mu)^\top (A\Sigma_p A^\top)^{-1} (A\Delta\mu)
&= \Delta\mu^\top A^\top A^{-\top}\Sigma_p^{-1}A^{-1}A\,\Delta\mu
 = \Delta\mu^\top \Sigma_p^{-1}\Delta\mu, \\[2pt]
\ln|A\Sigma_p A^\top| - \ln|A\Sigma_q A^\top|
&= \big(\ln|\Sigma_p| + 2\ln|A|\big) - \big(\ln|\Sigma_q| + 2\ln|A|\big)
 = \ln|\Sigma_p| - \ln|\Sigma_q|,
\end{align*}
where $\Delta\mu = \mu_q - \mu_p$ and the trace step uses the cyclic property. \hfill\qed
\end{proof}

\begin{lemma}[Invariance of the CEB objective and the MNI point]\label{lem:mi-invariance}
For any bijective measurable map $\phi$,
$I(X;\phi(Z)|Y) = I(X;Z|Y)$ and $I(Y;\phi(Z)) = I(Y;Z)$.
Consequently the CEB objective (Eq.~5) is invariant under bijections of $Z$, and
$Z$ attains the MNI point (Eq.~1) if and only if $\phi(Z)$ does.
\end{lemma}

\begin{proof}
Mutual information is invariant under bijective transformations of either argument
(data processing inequality applied in both directions). \hfill\qed
\end{proof}

\begin{proposition}[The FGIB side path is a CEB bottleneck on the main path]\label{prop:fgib-ceb}
Let $z = d$ and let the learned head $A$ be invertible (unconstrained). The \emph{fixed}
prior $r(z|y) = \mathcal{N}(\mu_p(y), \Sigma_p(y))$ --- for FGIB, the orthonormal anchor
table $\mu_p(y) = a\,v_y$, $\Sigma_p = I$ --- is, by construction, the image under $A$
of the canonical prior
$r'(\tilde{z}|y) = \mathcal{N}(A^{-1}\mu_p(y), A^{-1}\Sigma_p(y) A^{-\top})$
(since $A A^{-1} = I$), and the posterior $q(z|x) = \mathcal{N}(A\,h(x), \Sigma_q(x))$
is the image of
$q'(\tilde{z}|x) = \mathcal{N}(h(x), A^{-1}\Sigma_q(x) A^{-\top})$.
The prior is never transformed by training: $A^{-1}$ only reads the fixed anchor
geometry in the coordinates in which the posterior is centered at $h$.
By KL invariance under the common invertible map $A^{-1}$
(Lemma~\ref{lem:kl-invariance}), pointwise for every $(x, y)$,
\begin{equation}
\operatorname{KL}\big(q(z|x)\,\|\,r(z|y)\big)
= \operatorname{KL}\big(q'(\tilde{z}|x)\,\|\,r'(\tilde{z}|y)\big),
\label{eq:fgib-ceb}
\end{equation}
That is, the side-path regularizer is pointwise identical to a CEB bottleneck applied to
a noisy version of the main-path representation,
$\tilde{z} = h + \tilde{\epsilon}$ with
$\tilde{\epsilon} \sim \mathcal{N}(0, A^{-1}\Sigma_q A^{-\top})$,
using the canonical prior $r'$ as the backward encoder. For FGIB the canonical backward
encoder is the frame $\mathcal{N}(a\,A^{-1} v_y, A^{-1}A^{-\top})$, whose anchors
$A^{-1}v_y$ are orthonormal at scale $a$ in the induced metric $A^{\top}A$
($\langle A^{-1}v_i, A^{-1}v_j\rangle_{A^{\top}A} = \delta_{ij}$): the orthogonality
enters through the fixed frame $\{v_y\}$ and the equilibrium it induces, not through a
constraint on $A$.
Hence the regularizer's gradient with respect to $h$ (through the shared encoder)
coincides, up to the fixed linear map $A^{\top}$, with the gradient that a CEB model
applies to its representation mean.
\end{proposition}

\begin{proof}
The map $z \mapsto \tilde{z} = A^{-1}z$ is a bijection, and
$q' = (A^{-1})_\# q$, $r' = (A^{-1})_\# r$:
$A^{-1}A\,h = h$ and $A^{-1}\Sigma_q A^{-\top}$ give the canonical posterior, and
$A^{-1}\mu_p$, $A^{-1}\Sigma_p A^{-\top}$ give the canonical prior, with
$A\,(A^{-1}\mu_p) = \mu_p$ and $A\,(A^{-1}\Sigma_p A^{-\top})\,A^{\top} = \Sigma_p$, so
$r$ is exactly the image of $r'$ under $A$.
Equation~\eqref{eq:fgib-ceb} is Lemma~\ref{lem:kl-invariance} applied to $A^{-1}$.
Sampling $\tilde{z} \sim q'$ is equivalent to
$\tilde{z} = h + \tilde{\epsilon}$ with the stated noise.
Combined with the variational bound of Fischer (2020, Eqs.~9--11), the regularizer
upper-bounds the residual information of the noisy main-path representation:
$I(X;\tilde{Z}|Y) \le \mathbb{E}[\operatorname{KL}(q'(\tilde{z}|x)\,\|\,r'(\tilde{z}|y))]$.
\hfill\qed
\end{proof}

\begin{remark}[Family closure]\label{rem:closure}
Lemma~\ref{lem:kl-invariance} holds for the underlying Gaussian distributions. Within a
restricted variational family the equality of Proposition~\ref{prop:fgib-ceb} is exact
only if the family is closed under the inverse map $A^{-1}$: the full-covariance family
is closed (and is the parameterization used by Fischer 2020); the diagonal family is not
--- for diagonal $q, r$ the transformed pair is diagonal only if $A^{-1}$ preserves
diagonality (e.g., permutations and sign flips) and approximate otherwise. The diagonal
implementation of the codebase is therefore covered at the distribution level, with the
family approximation confined to the noise covariance $A^{-1}\Sigma_q A^{-\top}$
(fgib\_bound.tex, Remark on the unconstrained side head).
\end{remark}

\begin{remark}[Rank deficiency]\label{rem:rank}
If $z < d$ or $A$ is rank-deficient, only the projection of $h$ onto the image of $A$
receives gradient from the regularizer; the null-space component is unconstrained.
In the default configuration $z = d$ and a randomly initialized $A$ is almost surely invertible.
\end{remark}

\begin{remark}[Noise-free limit]\label{rem:limit}
The equivalence of Proposition~\ref{prop:fgib-ceb} should be understood at the level of the
variational objective, not as a mutual-information limit: as the side-path noise vanishes,
$I(X;Z|Y)$ diverges for continuous deterministic $H$ (the pathology of deterministic
networks discussed by Amjad \& Geiger and Kolchinsky et al.), and the variational KL bound
becomes loose. Keeping the side-path covariance finite --- without injecting noise into the
inference path --- is precisely what keeps the regularizer well-defined; this is the design
rationale of FGIB.
\end{remark}

% 若论文已定义 lemma/proposition/remark 环境，删除下面三行 \newtheorem

\newtheorem{lemma}{Lemma}
\newtheorem{proposition}{Proposition}
\newtheorem{remark}{Remark}

\textbf{Notation.}
Let $X$ be the input and $Y$ the label.
The main (deterministic) path of FGIB computes the representation
$H = f_{\mathrm{enc}}(X) \in \mathbb{R}^{d}$, and the classifier consumes $H$ directly.
The side path parameterizes $q(z|x) = \mathcal{N}(\mu_q(x), \Sigma_q(x))$ with linear mean
head $\mu_q(x) = A\,h(x)$ for $A \in \mathbb{R}^{z \times d}$, unconstrained ($z = d$, so
$A$ is a.s.\ invertible; the bias of the implementation's Linear head is omitted by the
same convention); $Z$ never feeds the classifier.
The label-conditional prior (backward encoder) is $r(z|y) = \mathcal{N}(\mu_p(y), \Sigma_p(y))$;
for FGIB it is the fixed orthonormal anchor table $\mathcal{N}(a v_y, I)$. The
orthogonality of the method lies in this fixed frame and in the equilibrium it induces
--- the KL drives the class-conditional means $\mu_y = A\,\mathbb{E}[h|y]$ toward the
orthonormal anchors $a v_y$, so the transformed class representations become mutually
orthogonal (the aligned family of fgib\_degeneration.tex) --- while $A$ itself remains
an unconstrained linear map.
The side-path regularizer is $\operatorname{KL}(q(z|x)\,\|\,r(z|y))$, added to the
cross-entropy with weight $\beta$, with gradients reaching the main path only through the
shared encoder (the classification gradient does not enter the side heads, and vice versa).
Following Fischer (2020), we write the residual information $I(X;Z|Y)$ (Eq.~4), the CEB
objective $\min_Z I(X;Z|Y) - \gamma I(Y;Z)$ (Eq.~5), the MNI point
$I(X;Y) = I(X;Z) = I(Y;Z)$ (Eq.~1), and the variational upper bound
$I(X;Z|Y) \le \mathbb{E}[\operatorname{KL}(q(z|x)\,\|\,r(z|y))]$ (Eqs.~9--11).

\textbf{Motivation.}
Stochastic bottlenecks such as VIB and CEB are \emph{stochastic at training time but
deterministic at evaluation time}: the classifier is trained on sampled codes $z \sim q(z|x)$
yet evaluated on the mean $\mu_q(x)$ (Fischer 2020, Appendix A.1), so its training and
evaluation distributions differ. In the same models the cross-entropy gradient and the KL
gradient act on the same heads: the classification term can exploit the variance (by Jensen,
$\mathbb{E}_z[\mathrm{CE}(c(z), y)] \ge \mathrm{CE}(c(\mu_q), y)$, so loosening the noise
relaxes the upper bound), while the KL term pushes the posterior back toward the prior.
Finally, the regularizer constrains the noisy code $z$, whereas inference consumes the
noise-free mean --- the object being regularized is not the object being used.
FGIB resolves all three mismatches by design: the classifier consumes the deterministic
representation $h$ at both training and evaluation time, the two objectives act on disjoint
parameter sets (the classification gradient does not enter the side heads, and the KL
gradient reaches the main path only through the shared encoder), and, by
Proposition~\ref{prop:fgib-ceb} below, the side-path regularizer is exactly a CEB bottleneck
applied to the representation that inference actually uses. Confining the stochasticity to a
side path also keeps the variational objective well-defined: for continuous deterministic
representations the residual information diverges (Amjad \& Geiger; Kolchinsky et al.),
so a purely deterministic bottleneck admits no finite variational bound, while the finite
side-path covariance of FGIB does (Remark~\ref{rem:limit}). Whether the robustness benefits
of the stochastic bottleneck survive the deterministic inference path is the empirical
question studied in the experiments.

\begin{lemma}[Linear invariance of the Gaussian KL]\label{lem:kl-invariance}
Let $q = \mathcal{N}(\mu_q, \Sigma_q)$ and $p = \mathcal{N}(\mu_p, \Sigma_p)$ on $\mathbb{R}^{z}$,
and let $A \in \mathbb{R}^{z \times z}$ be invertible. Define the common pushforwards
$q' = A_\# q = \mathcal{N}(A\mu_q, A\Sigma_q A^\top)$ and $p' = A_\# p$. Then
$\operatorname{KL}(q'\,\|\,p') = \operatorname{KL}(q\,\|\,p)$.
\end{lemma}

\begin{proof}
With $d = z$, the closed form reads
$\operatorname{KL}(q\,\|\,p)
= \tfrac{1}{2}\big[\operatorname{tr}(\Sigma_p^{-1}\Sigma_q)
+ (\mu_q - \mu_p)^\top \Sigma_p^{-1}(\mu_q - \mu_p)
- d + \ln|\Sigma_p| - \ln|\Sigma_q|\big]$.
Under the common map $A$ the three non-constant terms are unchanged:
\begin{align*}
\operatorname{tr}\big((A\Sigma_p A^\top)^{-1}(A\Sigma_q A^\top)\big)
&= \operatorname{tr}(A^{-\top}\Sigma_p^{-1}A^{-1} A\,\Sigma_q A^\top)
 = \operatorname{tr}(A^{-\top}\Sigma_p^{-1}\Sigma_q A^\top) \\
&= \operatorname{tr}(\Sigma_p^{-1}\Sigma_q A^\top A^{-\top})
 = \operatorname{tr}(\Sigma_p^{-1}\Sigma_q), \\[2pt]
(A\Delta\mu)^\top (A\Sigma_p A^\top)^{-1} (A\Delta\mu)
&= \Delta\mu^\top A^\top A^{-\top}\Sigma_p^{-1}A^{-1}A\,\Delta\mu
 = \Delta\mu^\top \Sigma_p^{-1}\Delta\mu, \\[2pt]
\ln|A\Sigma_p A^\top| - \ln|A\Sigma_q A^\top|
&= \big(\ln|\Sigma_p| + 2\ln|A|\big) - \big(\ln|\Sigma_q| + 2\ln|A|\big)
 = \ln|\Sigma_p| - \ln|\Sigma_q|,
\end{align*}
where $\Delta\mu = \mu_q - \mu_p$ and the trace step uses the cyclic property. \hfill\qed
\end{proof}

\begin{lemma}[Invariance of the CEB objective and the MNI point]\label{lem:mi-invariance}
For any bijective measurable map $\phi$,
$I(X;\phi(Z)|Y) = I(X;Z|Y)$ and $I(Y;\phi(Z)) = I(Y;Z)$.
Consequently the CEB objective (Eq.~5) is invariant under bijections of $Z$, and
$Z$ attains the MNI point (Eq.~1) if and only if $\phi(Z)$ does.
\end{lemma}

\begin{proof}
Mutual information is invariant under bijective transformations of either argument
(data processing inequality applied in both directions). \hfill\qed
\end{proof}

\begin{proposition}[The FGIB side path is a CEB bottleneck on the main path]\label{prop:fgib-ceb}
Let $z = d$ and let the learned head $A$ be invertible (unconstrained). The \emph{fixed}
prior $r(z|y) = \mathcal{N}(\mu_p(y), \Sigma_p(y))$ --- for FGIB, the orthonormal anchor
table $\mu_p(y) = a\,v_y$, $\Sigma_p = I$ --- is, by construction, the image under $A$
of the canonical prior
$r'(\tilde{z}|y) = \mathcal{N}(A^{-1}\mu_p(y), A^{-1}\Sigma_p(y) A^{-\top})$
(since $A A^{-1} = I$), and the posterior $q(z|x) = \mathcal{N}(A\,h(x), \Sigma_q(x))$
is the image of
$q'(\tilde{z}|x) = \mathcal{N}(h(x), A^{-1}\Sigma_q(x) A^{-\top})$.
The prior is never transformed by training: $A^{-1}$ only reads the fixed anchor
geometry in the coordinates in which the posterior is centered at $h$.
By KL invariance under the common invertible map $A^{-1}$
(Lemma~\ref{lem:kl-invariance}), pointwise for every $(x, y)$,
\begin{equation}
\operatorname{KL}\big(q(z|x)\,\|\,r(z|y)\big)
= \operatorname{KL}\big(q'(\tilde{z}|x)\,\|\,r'(\tilde{z}|y)\big),
\label{eq:fgib-ceb}
\end{equation}
That is, the side-path regularizer is pointwise identical to a CEB bottleneck applied to
a noisy version of the main-path representation,
$\tilde{z} = h + \tilde{\epsilon}$ with
$\tilde{\epsilon} \sim \mathcal{N}(0, A^{-1}\Sigma_q A^{-\top})$,
using the canonical prior $r'$ as the backward encoder. For FGIB the canonical backward
encoder is the frame $\mathcal{N}(a\,A^{-1} v_y, A^{-1}A^{-\top})$, whose anchors
$A^{-1}v_y$ are orthonormal at scale $a$ in the induced metric $A^{\top}A$
($\langle A^{-1}v_i, A^{-1}v_j\rangle_{A^{\top}A} = \delta_{ij}$): the orthogonality
enters through the fixed frame $\{v_y\}$ and the equilibrium it induces, not through a
constraint on $A$.
Hence the regularizer's gradient with respect to $h$ (through the shared encoder)
coincides, up to the fixed linear map $A^{\top}$, with the gradient that a CEB model
applies to its representation mean.
\end{proposition}

\begin{proof}
The map $z \mapsto \tilde{z} = A^{-1}z$ is a bijection, and
$q' = (A^{-1})_\# q$, $r' = (A^{-1})_\# r$:
$A^{-1}A\,h = h$ and $A^{-1}\Sigma_q A^{-\top}$ give the canonical posterior, and
$A^{-1}\mu_p$, $A^{-1}\Sigma_p A^{-\top}$ give the canonical prior, with
$A\,(A^{-1}\mu_p) = \mu_p$ and $A\,(A^{-1}\Sigma_p A^{-\top})\,A^{\top} = \Sigma_p$, so
$r$ is exactly the image of $r'$ under $A$.
Equation~\eqref{eq:fgib-ceb} is Lemma~\ref{lem:kl-invariance} applied to $A^{-1}$.
Sampling $\tilde{z} \sim q'$ is equivalent to
$\tilde{z} = h + \tilde{\epsilon}$ with the stated noise.
Combined with the variational bound of Fischer (2020, Eqs.~9--11), the regularizer
upper-bounds the residual information of the noisy main-path representation:
$I(X;\tilde{Z}|Y) \le \mathbb{E}[\operatorname{KL}(q'(\tilde{z}|x)\,\|\,r'(\tilde{z}|y))]$.
\hfill\qed
\end{proof}

\begin{remark}[Family closure]\label{rem:closure}
Lemma~\ref{lem:kl-invariance} holds for the underlying Gaussian distributions. Within a
restricted variational family the equality of Proposition~\ref{prop:fgib-ceb} is exact
only if the family is closed under the inverse map $A^{-1}$: the full-covariance family
is closed (and is the parameterization used by Fischer 2020); the diagonal family is not
--- for diagonal $q, r$ the transformed pair is diagonal only if $A^{-1}$ preserves
diagonality (e.g., permutations and sign flips) and approximate otherwise. The diagonal
implementation of the codebase is therefore covered at the distribution level, with the
family approximation confined to the noise covariance $A^{-1}\Sigma_q A^{-\top}$
(fgib\_bound.tex, Remark on the unconstrained side head).
\end{remark}

\begin{remark}[Rank deficiency]\label{rem:rank}
If $z < d$ or $A$ is rank-deficient, only the projection of $h$ onto the image of $A$
receives gradient from the regularizer; the null-space component is unconstrained.
In the default configuration $z = d$ and a randomly initialized $A$ is almost surely invertible.
\end{remark}

\begin{remark}[Noise-free limit]\label{rem:limit}
The equivalence of Proposition~\ref{prop:fgib-ceb} should be understood at the level of the
variational objective, not as a mutual-information limit: as the side-path noise vanishes,
$I(X;Z|Y)$ diverges for continuous deterministic $H$ (the pathology of deterministic
networks discussed by Amjad \& Geiger and Kolchinsky et al.), and the variational KL bound
becomes loose. Keeping the side-path covariance finite --- without injecting noise into the
inference path --- is precisely what keeps the regularizer well-defined; this is the design
rationale of FGIB.
\end{remark}

% 若论文已定义 lemma/proposition/remark 环境，删除下面三行 \newtheorem

\newtheorem{lemma}{Lemma}
\newtheorem{proposition}{Proposition}
\newtheorem{remark}{Remark}

\textbf{Notation.}
Let $X$ be the input and $Y$ the label.
The main (deterministic) path of FGIB computes the representation
$H = f_{\mathrm{enc}}(X) \in \mathbb{R}^{d}$, and the classifier consumes $H$ directly.
The side path parameterizes $q(z|x) = \mathcal{N}(\mu_q(x), \Sigma_q(x))$ with linear mean
head $\mu_q(x) = A\,h(x)$ for $A \in \mathbb{R}^{z \times d}$, unconstrained ($z = d$, so
$A$ is a.s.\ invertible; the bias of the implementation's Linear head is omitted by the
same convention); $Z$ never feeds the classifier.
The label-conditional prior (backward encoder) is $r(z|y) = \mathcal{N}(\mu_p(y), \Sigma_p(y))$;
for FGIB it is the fixed orthonormal anchor table $\mathcal{N}(a v_y, I)$. The
orthogonality of the method lies in this fixed frame and in the equilibrium it induces
--- the KL drives the class-conditional means $\mu_y = A\,\mathbb{E}[h|y]$ toward the
orthonormal anchors $a v_y$, so the transformed class representations become mutually
orthogonal (the aligned family of fgib\_degeneration.tex) --- while $A$ itself remains
an unconstrained linear map.
The side-path regularizer is $\operatorname{KL}(q(z|x)\,\|\,r(z|y))$, added to the
cross-entropy with weight $\beta$, with gradients reaching the main path only through the
shared encoder (the classification gradient does not enter the side heads, and vice versa).
Following Fischer (2020), we write the residual information $I(X;Z|Y)$ (Eq.~4), the CEB
objective $\min_Z I(X;Z|Y) - \gamma I(Y;Z)$ (Eq.~5), the MNI point
$I(X;Y) = I(X;Z) = I(Y;Z)$ (Eq.~1), and the variational upper bound
$I(X;Z|Y) \le \mathbb{E}[\operatorname{KL}(q(z|x)\,\|\,r(z|y))]$ (Eqs.~9--11).

\textbf{Motivation.}
Stochastic bottlenecks such as VIB and CEB are \emph{stochastic at training time but
deterministic at evaluation time}: the classifier is trained on sampled codes $z \sim q(z|x)$
yet evaluated on the mean $\mu_q(x)$ (Fischer 2020, Appendix A.1), so its training and
evaluation distributions differ. In the same models the cross-entropy gradient and the KL
gradient act on the same heads: the classification term can exploit the variance (by Jensen,
$\mathbb{E}_z[\mathrm{CE}(c(z), y)] \ge \mathrm{CE}(c(\mu_q), y)$, so loosening the noise
relaxes the upper bound), while the KL term pushes the posterior back toward the prior.
Finally, the regularizer constrains the noisy code $z$, whereas inference consumes the
noise-free mean --- the object being regularized is not the object being used.
FGIB resolves all three mismatches by design: the classifier consumes the deterministic
representation $h$ at both training and evaluation time, the two objectives act on disjoint
parameter sets (the classification gradient does not enter the side heads, and the KL
gradient reaches the main path only through the shared encoder), and, by
Proposition~\ref{prop:fgib-ceb} below, the side-path regularizer is exactly a CEB bottleneck
applied to the representation that inference actually uses. Confining the stochasticity to a
side path also keeps the variational objective well-defined: for continuous deterministic
representations the residual information diverges (Amjad \& Geiger; Kolchinsky et al.),
so a purely deterministic bottleneck admits no finite variational bound, while the finite
side-path covariance of FGIB does (Remark~\ref{rem:limit}). Whether the robustness benefits
of the stochastic bottleneck survive the deterministic inference path is the empirical
question studied in the experiments.

\begin{lemma}[Linear invariance of the Gaussian KL]\label{lem:kl-invariance}
Let $q = \mathcal{N}(\mu_q, \Sigma_q)$ and $p = \mathcal{N}(\mu_p, \Sigma_p)$ on $\mathbb{R}^{z}$,
and let $A \in \mathbb{R}^{z \times z}$ be invertible. Define the common pushforwards
$q' = A_\# q = \mathcal{N}(A\mu_q, A\Sigma_q A^\top)$ and $p' = A_\# p$. Then
$\operatorname{KL}(q'\,\|\,p') = \operatorname{KL}(q\,\|\,p)$.
\end{lemma}

\begin{proof}
With $d = z$, the closed form reads
$\operatorname{KL}(q\,\|\,p)
= \tfrac{1}{2}\big[\operatorname{tr}(\Sigma_p^{-1}\Sigma_q)
+ (\mu_q - \mu_p)^\top \Sigma_p^{-1}(\mu_q - \mu_p)
- d + \ln|\Sigma_p| - \ln|\Sigma_q|\big]$.
Under the common map $A$ the three non-constant terms are unchanged:
\begin{align*}
\operatorname{tr}\big((A\Sigma_p A^\top)^{-1}(A\Sigma_q A^\top)\big)
&= \operatorname{tr}(A^{-\top}\Sigma_p^{-1}A^{-1} A\,\Sigma_q A^\top)
 = \operatorname{tr}(A^{-\top}\Sigma_p^{-1}\Sigma_q A^\top) \\
&= \operatorname{tr}(\Sigma_p^{-1}\Sigma_q A^\top A^{-\top})
 = \operatorname{tr}(\Sigma_p^{-1}\Sigma_q), \\[2pt]
(A\Delta\mu)^\top (A\Sigma_p A^\top)^{-1} (A\Delta\mu)
&= \Delta\mu^\top A^\top A^{-\top}\Sigma_p^{-1}A^{-1}A\,\Delta\mu
 = \Delta\mu^\top \Sigma_p^{-1}\Delta\mu, \\[2pt]
\ln|A\Sigma_p A^\top| - \ln|A\Sigma_q A^\top|
&= \big(\ln|\Sigma_p| + 2\ln|A|\big) - \big(\ln|\Sigma_q| + 2\ln|A|\big)
 = \ln|\Sigma_p| - \ln|\Sigma_q|,
\end{align*}
where $\Delta\mu = \mu_q - \mu_p$ and the trace step uses the cyclic property. \hfill\qed
\end{proof}

\begin{lemma}[Invariance of the CEB objective and the MNI point]\label{lem:mi-invariance}
For any bijective measurable map $\phi$,
$I(X;\phi(Z)|Y) = I(X;Z|Y)$ and $I(Y;\phi(Z)) = I(Y;Z)$.
Consequently the CEB objective (Eq.~5) is invariant under bijections of $Z$, and
$Z$ attains the MNI point (Eq.~1) if and only if $\phi(Z)$ does.
\end{lemma}

\begin{proof}
Mutual information is invariant under bijective transformations of either argument
(data processing inequality applied in both directions). \hfill\qed
\end{proof}

\begin{proposition}[The FGIB side path is a CEB bottleneck on the main path]\label{prop:fgib-ceb}
Let $z = d$ and let the learned head $A$ be invertible (unconstrained). The \emph{fixed}
prior $r(z|y) = \mathcal{N}(\mu_p(y), \Sigma_p(y))$ --- for FGIB, the orthonormal anchor
table $\mu_p(y) = a\,v_y$, $\Sigma_p = I$ --- is, by construction, the image under $A$
of the canonical prior
$r'(\tilde{z}|y) = \mathcal{N}(A^{-1}\mu_p(y), A^{-1}\Sigma_p(y) A^{-\top})$
(since $A A^{-1} = I$), and the posterior $q(z|x) = \mathcal{N}(A\,h(x), \Sigma_q(x))$
is the image of
$q'(\tilde{z}|x) = \mathcal{N}(h(x), A^{-1}\Sigma_q(x) A^{-\top})$.
The prior is never transformed by training: $A^{-1}$ only reads the fixed anchor
geometry in the coordinates in which the posterior is centered at $h$.
By KL invariance under the common invertible map $A^{-1}$
(Lemma~\ref{lem:kl-invariance}), pointwise for every $(x, y)$,
\begin{equation}
\operatorname{KL}\big(q(z|x)\,\|\,r(z|y)\big)
= \operatorname{KL}\big(q'(\tilde{z}|x)\,\|\,r'(\tilde{z}|y)\big),
\label{eq:fgib-ceb}
\end{equation}
That is, the side-path regularizer is pointwise identical to a CEB bottleneck applied to
a noisy version of the main-path representation,
$\tilde{z} = h + \tilde{\epsilon}$ with
$\tilde{\epsilon} \sim \mathcal{N}(0, A^{-1}\Sigma_q A^{-\top})$,
using the canonical prior $r'$ as the backward encoder. For FGIB the canonical backward
encoder is the frame $\mathcal{N}(a\,A^{-1} v_y, A^{-1}A^{-\top})$, whose anchors
$A^{-1}v_y$ are orthonormal at scale $a$ in the induced metric $A^{\top}A$
($\langle A^{-1}v_i, A^{-1}v_j\rangle_{A^{\top}A} = \delta_{ij}$): the orthogonality
enters through the fixed frame $\{v_y\}$ and the equilibrium it induces, not through a
constraint on $A$.
Hence the regularizer's gradient with respect to $h$ (through the shared encoder)
coincides, up to the fixed linear map $A^{\top}$, with the gradient that a CEB model
applies to its representation mean.
\end{proposition}

\begin{proof}
The map $z \mapsto \tilde{z} = A^{-1}z$ is a bijection, and
$q' = (A^{-1})_\# q$, $r' = (A^{-1})_\# r$:
$A^{-1}A\,h = h$ and $A^{-1}\Sigma_q A^{-\top}$ give the canonical posterior, and
$A^{-1}\mu_p$, $A^{-1}\Sigma_p A^{-\top}$ give the canonical prior, with
$A\,(A^{-1}\mu_p) = \mu_p$ and $A\,(A^{-1}\Sigma_p A^{-\top})\,A^{\top} = \Sigma_p$, so
$r$ is exactly the image of $r'$ under $A$.
Equation~\eqref{eq:fgib-ceb} is Lemma~\ref{lem:kl-invariance} applied to $A^{-1}$.
Sampling $\tilde{z} \sim q'$ is equivalent to
$\tilde{z} = h + \tilde{\epsilon}$ with the stated noise.
Combined with the variational bound of Fischer (2020, Eqs.~9--11), the regularizer
upper-bounds the residual information of the noisy main-path representation:
$I(X;\tilde{Z}|Y) \le \mathbb{E}[\operatorname{KL}(q'(\tilde{z}|x)\,\|\,r'(\tilde{z}|y))]$.
\hfill\qed
\end{proof}

\begin{remark}[Family closure]\label{rem:closure}
Lemma~\ref{lem:kl-invariance} holds for the underlying Gaussian distributions. Within a
restricted variational family the equality of Proposition~\ref{prop:fgib-ceb} is exact
only if the family is closed under the inverse map $A^{-1}$: the full-covariance family
is closed (and is the parameterization used by Fischer 2020); the diagonal family is not
--- for diagonal $q, r$ the transformed pair is diagonal only if $A^{-1}$ preserves
diagonality (e.g., permutations and sign flips) and approximate otherwise. The diagonal
implementation of the codebase is therefore covered at the distribution level, with the
family approximation confined to the noise covariance $A^{-1}\Sigma_q A^{-\top}$
(fgib\_bound.tex, Remark on the unconstrained side head).
\end{remark}

\begin{remark}[Rank deficiency]\label{rem:rank}
If $z < d$ or $A$ is rank-deficient, only the projection of $h$ onto the image of $A$
receives gradient from the regularizer; the null-space component is unconstrained.
In the default configuration $z = d$ and a randomly initialized $A$ is almost surely invertible.
\end{remark}

\begin{remark}[Noise-free limit]\label{rem:limit}
The equivalence of Proposition~\ref{prop:fgib-ceb} should be understood at the level of the
variational objective, not as a mutual-information limit: as the side-path noise vanishes,
$I(X;Z|Y)$ diverges for continuous deterministic $H$ (the pathology of deterministic
networks discussed by Amjad \& Geiger and Kolchinsky et al.), and the variational KL bound
becomes loose. Keeping the side-path covariance finite --- without injecting noise into the
inference path --- is precisely what keeps the regularizer well-defined; this is the design
rationale of FGIB.
\end{remark}
